%============================================================
% CONCLUSION
%============================================================

\section{Conclusion}
\label{sec:conclusion}

This study investigated the distribution of baseline headache
frequency using Poisson finite mixture models. The analysis began
with a standard Poisson model and then considered finite mixtures to
account for the substantial heterogeneity observed among patients.

The exploratory analysis showed clear over-dispersion, with a
variance-to-mean ratio of 2.78. The NPMLE analysis provided further
evidence that the observed population was unlikely to be adequately
represented by a single homogeneous Poisson distribution.

Finite mixture models with two, three, four, and five components were
compared. Although the Bayesian Information Criterion favored the
two-component model, the Akaike Information Criterion favored the
three-component model. The three-component solution was therefore
selected for detailed interpretation because it provided the lowest
AIC and produced a clear and interpretable separation of the
population into low-, intermediate-, and high-frequency headache
groups.

The estimated component means were approximately 9.61, 13.80, and
22.95 headache days, demonstrating substantial differences in
underlying headache-frequency rates between the latent groups.
Results obtained using CAMAN and FlexMix showed broadly similar
component means, providing additional support for the identified
heterogeneity.

The analysis also demonstrated an important limitation of the fitted
Poisson mixture. The model substantially underestimated the number of
patients reporting the maximum value of 28 headache days. This
suggests that the bounded nature of the four-week observation period
should be considered in future modelling.

Including age as a covariate showed that its association with
headache frequency was not uniform across the latent groups. Age was
statistically significant in one component, where a one-year increase
in age was associated with an estimated 2.4\% decrease in the
expected headache frequency, while no significant age effects were
identified in the other two components.

Overall, the findings demonstrate that Poisson finite mixture models
provide a useful framework for investigating unobserved heterogeneity
in headache-frequency data. The analysis also illustrates the
importance of combining exploratory methods, model-selection
criteria, graphical diagnostics, and substantive interpretation when
choosing and evaluating a finite mixture model.

Future work could consider models that explicitly account for the
upper boundary of the headache-frequency outcome and could
incorporate additional clinical or demographic covariates to further
explain the observed heterogeneity.